\chapter{Pruebas}

\section{Estrategia de validación}

Un sistema basado en agentes no se valida solo comprobando tipos de retorno. Hay
que verificar coordinación entre etapas, persistencia, revisión humana,
recuperación de memoria, coste y calidad de los artefactos científicos.

\begin{table}[H]
\begin{center}
\small
\begin{tabular}{|p{3cm}|p{9.5cm}|} \hline
\textbf{Nivel} & \textbf{Validación} \\ \hline
Estructura & JSON, modelos Pydantic y contratos de herramientas. \\ \hline
Pipeline & Estados del Router, persistencia y reanudación. \\ \hline
Código generado & Importación del modelo y tests sobre \texttt{decide}, \texttt{update} y \texttt{get\_state}. \\ \hline
Memoria & Escritura tras revisión, recuperación, deduplicación y validez temporal. \\ \hline
Integración & Carga dinámica del modelo desde la fase de simulación. \\ \hline
Evals & Suites reproducibles para medir pipeline, memoria, coste, latencia y calidad. \\ \hline
\end{tabular}
\caption{Niveles de validación.}
\end{center}
\end{table}

La evidencia versionada incluye tests unitarios, suites YAML, ejecuciones de
ejemplo, bundles de evals y reportes de juez. Esto permite revisar el resultado
sin depender de una conversación concreta con un modelo.
Las Secciones~\ref{ap:validacion-extendida}
y~\ref{ap:limitaciones-extendidas} conservan la explicación completa de CI,
evals, corpus cerrado, reservas del juez y amenazas a la validez.

\begin{table}[H]
\begin{center}
\small
\begin{tabular}{|p{4.1cm}|p{2.1cm}|p{6.1cm}|} \hline
\textbf{Evidencia} & \textbf{Cantidad} & \textbf{Qué demuestra} \\ \hline
Informes del \texttt{sample-run} & 7 deep reports & Investigación estructurada por paradigma. \\ \hline
Formulaciones del \texttt{sample-run} & 7 documentos & Conversión de investigación en variables y reglas. \\ \hline
Specs del Reasoner & 24 specs JSON & Adaptación de formulaciones al entorno. \\ \hline
Builder del \texttt{sample-run} & 2 modelos + tests & Código Python autocontenido y validado. \\ \hline
Tests de fase 1 y shared & 131 archivos & Agentes, runtime, memoria, herramientas e infraestructura. \\ \hline
Suites YAML & 10 suites & Evals reproducibles de pipeline, memoria y recuperación. \\ \hline
Corpus cerrado & 2 casos, 6 PDFs & Ejecuciones completas con corpus limitado y juez externo. \\ \hline
\end{tabular}
\caption{Evidencias versionadas de validación.}
\end{center}
\end{table}

\subsection{Cierre de requisitos}

La tabla siguiente cierra los requisitos funcionales contra evidencia
versionada. El estado no significa certificación científica final. Indica que el
comportamiento exigido existe, se ejecuta y deja evidencia revisable.

\begin{table}[H]
\begin{center}
\scriptsize
\begin{tabular}{|p{1.5cm}|p{2.0cm}|p{6.2cm}|p{3.0cm}|} \hline
\textbf{Req.} & \textbf{Estado} & \textbf{Evidencia principal} & \textbf{Reserva} \\ \hline
RF-01 & Cumplido & CLI, servidor y suites de pipeline arrancan runs desde descripción. & Depende de proveedor LLM. \\ \hline
RF-02 & Cumplido & Classifier y evals de slugs miden reutilización canónica. & Hay alerta histórica de fragmentación. \\ \hline
RF-03 & Cumplido & \texttt{sample-run}, deep reports y corpus cerrado generan informes. & Calidad final requiere revisión. \\ \hline
RF-04 & Cumplido & Router y \texttt{FeedbackPort} bloquean antes de formalizar. & En evals se usa autoaprobación. \\ \hline
RF-05 & Cumplido & 7 documentos de formulación en \texttt{sample-run}. & Algunas variantes pueden ser descartadas. \\ \hline
RF-06 & Cumplido & Selección humana o autoaprobada antes de Reasoner. & La selección afecta cobertura final. \\ \hline
RF-07 & Cumplido & Reasoner usa \texttt{env\_spec.json} como fuente autoritativa. & Specs pobres limitan el resultado. \\ \hline
RF-08 & Cumplido & 24 specs JSON y salidas inválidas explícitas. & CASO1 incluye una spec inválida. \\ \hline
RF-09 & Cumplido & Modelos con \texttt{decide}, \texttt{update} y \texttt{get\_state}. & Puede haber problemas de registro. \\ \hline
RF-10 & Cumplido & Builder genera tests y ejecuta \texttt{pytest}. & Tests prueban contrato, no validez científica. \\ \hline
RF-11 & Cumplido & MinIO, PostgreSQL, bundles y modelos registrados. & Requiere servicios centrales activos. \\ \hline
RF-12 & Cumplido con reservas & MemoryAgent escribe tras aceptación y evals de memoria. & Calidad de extracción depende del artefacto. \\ \hline
RF-13 & Cumplido & \texttt{pipeline\_state.json}, \texttt{trace.jsonl} y Agrex. & Reanudación exige artefactos presentes. \\ \hline
\end{tabular}
\caption{Cierre de requisitos funcionales.}
\end{center}
\end{table}

\begin{table}[H]
\begin{center}
\scriptsize
\begin{tabular}{|p{1.45cm}|p{1.8cm}|p{6.2cm}|p{2.45cm}|} \hline
\textbf{Req.} & \textbf{Estado} & \textbf{Evidencia principal} & \textbf{Reserva} \\ \hline
RNF-01 & Cumplido & Artefactos, trazas, estado y reportes por run. & Auditoría depende de persistencia. \\ \hline
RNF-02 & Cumplido & Estado persistido y comprobación de artefactos. & No reanuda si faltan objetos. \\ \hline
RNF-03 & Cumplido & Agentes y subagentes separados por responsabilidad. & Añade coordinación. \\ \hline
RNF-04 & Cumplido & Fan-out por paradigmas, formulaciones y specs. & Coste crece con ramas. \\ \hline
RNF-05 & Cumplido & Contrato \texttt{DecisionModel} estable. & Nuevas memorias exigen adaptadores. \\ \hline
RNF-06 & Cumplido & LLMs, búsqueda, S3, bases de datos y fase 2. & Integra servicios externos. \\ \hline
RNF-07 & Cumplido & Modo degradado sin Neo4j, Qdrant o embeddings. & Sin memoria avanzada. \\ \hline
RNF-08 & Cumplido & Tests generados y suite \texttt{full-pipeline}. & Tests no sustituyen expertos. \\ \hline
RNF-09 & Cumplido & Presupuestos y costes en reportes de evals. & Suites online siguen siendo caras. \\ \hline
\end{tabular}
\caption{Cierre de requisitos no funcionales.}
\end{center}
\end{table}

\section{Evals y benchmarks}

Las evals se definen como suites YAML. Cada suite declara etapas, entradas,
aprobaciones automáticas, presupuesto máximo y aserciones esperadas. El reporte
incluye estado, tiempos, costes, crecimiento del grafo, llamadas a herramientas
y aserciones superadas o fallidas.

\begin{table}[H]
\begin{center}
\scriptsize
\begin{tabular}{|p{3.0cm}|p{3.1cm}|p{6.2cm}|} \hline
\textbf{Grupo} & \textbf{Suites} & \textbf{Qué valida} \\ \hline
Pipeline & Smoke, full pipeline & Coordinación de agentes, artefactos compatibles, modelos importables y \texttt{decide} funcional. \\ \hline
Memoria & Canonicalization, growth, slugs, retrieval & Reutilización de identidad científica, crecimiento controlado y recuperación de paradigmas previos. \\ \hline
Recuperación offline & Retrieval quality, temporal & Hechos recuperados en top-k, validez temporal y supersesión. \\ \hline
Corpus cerrado & CASO1 y CASO2 PDF corpus & Descubrimiento, formalización, adaptación y construcción usando solo papers suministrados. \\ \hline
\end{tabular}
\caption{Grupos de evals.}
\end{center}
\end{table}

\subsection{Resultados representativos}

\begin{table}[H]
\begin{center}
\small
\begin{tabular}{|p{2.8cm}|p{1.5cm}|p{2.5cm}|p{5.6cm}|} \hline
\textbf{Suite} & \textbf{Estado} & \textbf{Tiempo / coste} & \textbf{Evidencia} \\ \hline
Full pipeline & PASS & 2445.9 s / \$9.30 & Generó modelos, los importó y \texttt{decide} devolvió acciones válidas. \\ \hline
Retrieval quality & PASS & 22.7 s / \$0.08 & Cuatro consultas recuperaron hechos esperados en top-k. \\ \hline
Temporal correctness & PASS & 0.1 s / \$0.00 & Validó presencia, ausencia y cadena de supersesión. \\ \hline
Slug accuracy histórica & Alerta & 3646.2 s / \$5.08 & Detectó fragmentación de slugs y motivó restricciones de canonicalización. \\ \hline
\end{tabular}
\caption{Resultados representativos de evals internas.}
\end{center}
\end{table}

La fila de \texttt{slug-accuracy} se conserva como alerta histórica. No invalida
el estado final, pero muestra por qué fue necesario medir identidad científica y
no solo éxito técnico.

\subsection{Evaluación con corpus cerrado de papers}
\label{sec:evaluacion-corpus-cerrado}

La evaluación más exigente se hizo con 2 casos basados en PDFs. En cada caso,
el sistema solo podía acceder a los papers incluidos en el ZIP de entrada. Un
LLM juez revisó artefactos congelados con la misma rúbrica: fidelidad al corpus,
calidad de paradigmas, formulaciones, modelos, memoria, trazabilidad y
advertencias.

Esta prueba separa la validez técnica de la validez científica. La validez
técnica pregunta si el pipeline termina, registra artefactos, genera modelos y
ejecuta tests. La validez científica pregunta si los paradigmas y formulaciones
son fieles al corpus y si un experto podría defender el modelo. Las dos son
necesarias, pero no equivalentes.

\begin{table}[H]
\begin{center}
\small
\begin{tabular}{|p{3.1cm}|p{4.9cm}|p{4.5cm}|} \hline
\textbf{Lectura} & \textbf{Qué demuestra} & \textbf{Qué no demuestra} \\ \hline
PASS técnico & El sistema completa etapas, persiste artefactos y genera modelos ejecutables. & Que el modelo represente perfectamente el paradigma científico. \\ \hline
Juez externo & Un evaluador separado revisa fidelidad, taxonomía, trazabilidad y reservas. & Que la evaluación sustituya a un experto humano del dominio. \\ \hline
Aprobado con reservas & Los artefactos son útiles y auditables, pero tienen fallos localizados. & Que el resultado pueda aceptarse sin revisión posterior. \\ \hline
\end{tabular}
\caption{Diferencia entre validación técnica y científica.}
\end{center}
\end{table}

\begin{table}[H]
\begin{center}
\small
\begin{tabular}{|p{1.6cm}|p{3.1cm}|p{2.4cm}|p{5.1cm}|} \hline
\textbf{Caso} & \textbf{Resultado} & \textbf{Tiempo / coste} & \textbf{Evidencia generada} \\ \hline
CASO1 & Pipeline PASS; juez 72/100, aprobado con reservas. & 6152.0 s / \$33.71 & 3 papers, 6 paradigmas, 18 specs, 15 modelos y KG de 433 nodos. \\ \hline
CASO2 & Pipeline PASS; juez 78/100, aprobado con reservas. & 4961.8 s / \$24.47 & 3 papers, 4 paradigmas, 12 specs, 12 modelos y KG de 272 nodos. \\ \hline
\end{tabular}
\caption{Resultados de las evals con corpus cerrado.}
\end{center}
\end{table}

El resultado fue positivo, pero no perfecto. Esa lectura es importante: el
pipeline puede ser técnicamente correcto y aun así requerir revisión científica.

\begin{table}[H]
\begin{center}
\small
\begin{tabular}{|p{3.0cm}|p{2.1cm}|p{7.0cm}|} \hline
\textbf{Reserva} & \textbf{Severidad} & \textbf{Lectura} \\ \hline
Spec inválida en CASO1 & Alta & Reasoner detectó una formulación que no debía generar modelo aceptado. \\ \hline
Desajuste de registro en DDM & Alta & El bundle permitió localizar que el registro podía apuntar a clase o test distinto. \\ \hline
Errores de escritura en KG & Media & Reducen trazabilidad automática, pero no eliminan artefactos ni reportes. \\ \hline
Solapamiento taxonómico & Media & CASO2 mezcló mecanismos cercanos como homeostasis, HRL y active inference. \\ \hline
Duplicados de ficheros & Baja & Afectan higiene de revisión, no el contrato ejecutable del modelo. \\ \hline
\end{tabular}
\caption{Reservas detectadas por el juez externo.}
\end{center}
\end{table}

Por tanto, ``aprobado con reservas'' debe leerse como señal de utilidad, no
como certificación científica final. La eval reduce el riesgo de evaluar solo
que el sistema termina, porque obliga a revisar si los artefactos son fieles al
corpus y auditables por una persona.

Las reservas también orientan trabajo futuro. La spec inválida muestra que el
Reasoner debe poder rechazar una formulación en lugar de forzar código. El
desajuste de registro muestra que no basta con generar ficheros: hay que auditar
la ruta que se registra como modelo canónico. Los errores del grafo muestran que
MinIO y PostgreSQL deben conservar evidencia aunque falle una escritura de
Neo4j. El solapamiento taxonómico muestra que la revisión humana sigue siendo
necesaria en paradigmas cercanos.

\section{Pruebas técnicas}

El Builder genera tanto el modelo como sus pruebas. Después ejecuta:

\begin{verbatim}
uv run pytest
\end{verbatim}

Las pruebas comprueban que \texttt{decide} devuelve acciones válidas, que
\texttt{update} concentra la mutación de estado y que \texttt{get\_state}
expone datos serializables. Además, cada ejecución conserva
\texttt{pipeline\_state.json}, \texttt{trace.jsonl}, informes, specs, modelos y
tests para auditoría.

\section{Amenazas a la validez}

\begin{table}[H]
\begin{center}
\small
\begin{tabular}{|p{3cm}|p{4.9cm}|p{4.5cm}|} \hline
\textbf{Amenaza} & \textbf{Riesgo} & \textbf{Mitigación aplicada} \\ \hline
Validez científica & Un modelo puede pasar tests sin representar bien el paradigma. & Revisión humana y juez externo con rúbrica. \\ \hline
Validez interna & Los LLMs pueden variar con prompts, proveedor o contexto. & Suites versionadas, reportes JSON, costes y trazas. \\ \hline
Validez externa & Solo se ejecutaron 2 casos con corpus cerrado. & Se declaran como evidencia interna, no benchmark externo. \\ \hline
Validez de construcción & El juez LLM puede compartir sesgos con el sistema. & El juez revisa artefactos congelados y no participa en generación. \\ \hline
Reproducibilidad & Las ejecuciones completas dependen de servicios y presupuesto. & Bundles, snapshots y rutas versionadas. \\ \hline
\end{tabular}
\caption{Amenazas a la validez y mitigaciones.}
\end{center}
\end{table}
